\section{Related Work}\label{sec:related}

\paragraph{Decidability of TCTL satisfiability.}
Satisfiability of full TCTL is undecidable~\cite{AlurCourcoubetisDill1990,AlurHenzinger1994ReallyTemporal},
the source of undecidability being {punctual} timing constraints. La~Torre and
Napoli~\cite{fsttcs2000-LN,LaTorreNapoli2001TimedTreeAutomata} delimited the decidability frontier
by identifying the {non-punctual} fragment -- timing intervals may not be singular ($[a,a]$ is
forbidden), while strict and non-strict bounds are both allowed -- and decided it in \textsc{ExpTime}
by reducing satisfiability to emptiness of timed tree automata. We target exactly this fragment, with
matching complexity, but replace the global automata reduction by a goal-directed tableau: clock
valuations are kept symbolically as zones, only the demanded part of the state space is built, the
calculus filters out physically unrealisable (Zeno) runs explicitly, and a successful search returns
a witness model -- none of which the automata-theoretic decision procedure provides directly.

\paragraph{Bounded and parametric approaches.}
Penczek et al.~\cite{boundedTCTL} also address TCTL satisfiability, but from a bounded-synthesis
angle: satisfiability is reduced to the existence of a bounded parametric timed automaton whose
parameters are instantiated with the \textsc{Imitator} model checker. Termination stems from the
imposed structural bound rather than from a decision procedure, so the method is incomplete; its
fragment excludes nested modalities and permits punctual constraints, for which the problem stays
undecidable. We instead decide a decidable fragment, without nesting restrictions, completely.

\paragraph{Tableaux for branching-time logics.}
Tableaux are the standard proof-search decision procedures for untimed branching-time
logics~\cite{Clarke99,tableau_ctl}; we adapt the one-pass, context-based CTL tableau
of~\cite{tableau_ctl} to dense time. The closest quantitative predecessor is Emerson et
al.~\cite{EmersonEtAl1992QuantitativeTemporalReasoning}, a tableau for a CTL with {discrete}-time
bounds. Discreteness sidesteps the two phenomena that drive our development -- symbolic dense clock
valuations and time convergence -- so, while it demonstrates that tableaux extend to quantitative
branching time, it does not yield a procedure for the dense-time fragment considered here.

\paragraph{Model checking and \textsc{Uppaal}.}
Most real-time verification effort targets {model checking} -- TCTL against a fixed timed
automaton -- where \textsc{Uppaal}~\cite{BehrmannLarsenYi2004Uppaal} is the de-facto standard. It is
worth contrasting the {specification logic} the two settings use. \textsc{Uppaal}'s query
language is a deliberately {simplified}, flat fragment of TCTL: Boolean state properties under a
single, non-nested path modality ($\mathrm{AG}$, $\mathrm{EF}$, $\mathrm{AF}$, $\mathrm{EG}$) plus a
leads-to operator $\mathrm{AG}(\varphi\rightarrow\mathrm{AF}\,\psi)$,
with timing expressed through clocks in the {model} rather than through nested bounded
modalities in the formula. This keeps queries efficiently checkable but cannot express the nested,
time-bounded branching properties our calculus decides. The flat \textsc{Uppaal} fragment is a
syntactic restriction of our logic, so our procedure decides it as a special case; we use it later as
a practical reference point for satisfiability of realistic specifications.
